\section{Vanishing entropy}\label{sec:entropy}

Proposition~\ref{prop:attraction} gives pointwise attraction to
$Z_0\cup Z_\infty$.  We first prove that such attraction transfers zero
entropy from a compact invariant set to the whole space.  The argument
uses a uniform bound on the proportion of orbit blocks outside a
neighborhood of that set; it requires no uniform eventual-entry time.

The variational principle also yields this conclusion
\cite[Section~3]{Misiurewicz1976}: invariance and bounded convergence
force every invariant probability measure to be supported on the
attracting set. The direct proof below makes the exceptional-block
estimate explicit.

\begin{lemma}[Entropy localization]\label{lemma:localization}
Let $F:X\to X$ be a continuous map of a nonempty compact metric space
$(X,d)$, and let $Z\subseteq X$ be a nonempty compact set with
$F(Z)\subseteq Z$.  Suppose that
\[
  h_{\mathrm{top}}(F|_Z)=0
  \quad\text{and}\quad
  d(F^n x,Z)\longrightarrow 0
  \quad\text{for every }x\in X.
\]
Then $h_{\mathrm{top}}(F)=0$.
\end{lemma}

\begin{proof}
Write $s_X(n,\varepsilon)$ and $s_Z(n,\varepsilon)$ for the separated-set
cardinalities of $F$ and $F|_Z$, respectively.  Fix $\varepsilon>0$ and
an integer $m\geq1$, and define
\[
  d_m(x,y)=\max_{0\leq j<m}d(F^j x,F^j y).
\]
This metric induces the original topology of $X$.  Choose a largest
$(m,\varepsilon/4)$-separated subset of $Z$, of cardinality
$P=P(m)=s_Z(m,\varepsilon/4)$.  Its open $d_m$-balls of radius
$\varepsilon/2$, taken in $X$, have union $V$ containing $Z$: otherwise a
point of $Z$ at $d_m$-distance greater than $\varepsilon/4$ from every
center could be added to the separated set.  Also choose $Q=Q(m)$
open $d_m$-balls of radius $\varepsilon/2$ covering $X$.

We claim that for every $0<\tau<1$ there is an integer $J\geq1$ such
that
\begin{equation}\label{eq:sparse-entropy-blocks}
  \#\{0\leq i<r:F^{mi}x\notin V\}\leq \tau r+J
  \qquad(x\in X,\ r\geq1).
\end{equation}
Indeed, since $Z$ is compact and $V$ is an open neighborhood of $Z$,
the attraction hypothesis implies that $F^{mi}x\in V$ for all
sufficiently large $i$, with the threshold allowed to depend on $x$.
Consequently, for each $x$ there is an integer $k_x\geq1$ for which
\[
  \#\{0\leq i<k_x:F^{mi}x\notin V\}<\tau k_x.
\]
Let $I_x=\{0\leq i<k_x:F^{mi}x\in V\}$.  The open neighborhood
\[
  O_x=\bigcap_{i\in I_x}(F^{mi})^{-1}(V)
\]
of $x$ preserves all these visits: every $y\in O_x$ has fewer than
$\tau k_x$ indices outside $V$ among its first $k_x$ samples.
Choose a finite cover $O_{x_1},\ldots,O_{x_a}$ of $X$, and put
$J=\max_{1\leq b\leq a}k_{x_b}$.

Starting at $x$, choose a member of this cover containing the current
point and follow the corresponding number $k_{x_b}$ of iterates of
$F^m$.  Restart the same procedure at the next point, concatenating
these segments for as long as the selected segment fits within the
first $r$ samples.  The final remainder has length $\ell<J$.
The completed segments contribute at most $\tau(r-\ell)$ samples
outside $V$, and the remainder contributes at most $\ell$.  This
proves~\eqref{eq:sparse-entropy-blocks}, including when no segment
fits.  In particular, this compactness argument bounds a proportion
of exceptional samples without bounding their last occurrence.

Now encode the first $r$ blocks of length $m$ of an orbit.  Record
the subset of indices $i$ for which $F^{mi}x\notin V$.  At every other
index choose one of the $P$ balls covering $V$ that contains
$F^{mi}x$; at an exceptional index choose one of the $Q$ balls
covering $X$.  Choose the first available ball in a fixed ordering
of each cover.  By~\eqref{eq:sparse-entropy-blocks}, the number of
possible descriptions is at most
\begin{equation}\label{eq:entropy-description-count}
  2^r P^r Q^{\tau r+J}.
\end{equation}
Here $2^r$ bounds the number of subsets, and $P,Q\geq1$ allow the
displayed overcount even if $\tau r+J\geq r$.  Two points with the
same description lie in the same radius-$\varepsilon/2$ ball at each
block, so their distance is less than $\varepsilon$ at every time
$0,\ldots,mr-1$.  Thus an $(mr,\varepsilon)$-separated set has at most
the cardinality in~\eqref{eq:entropy-description-count}.

For arbitrary $n$, use $r=\lceil n/m\rceil$ and
$s_X(n,\varepsilon)\leq s_X(mr,\varepsilon)$.  Taking the orbit-length
limit with $m$ and $\tau$ fixed gives
\[
  \limsup_{n\to\infty}\frac{1}{n}\log s_X(n,\varepsilon)
  \leq \frac{\log 2+\log P(m)+\tau\log Q(m)}{m}.
\]
First let $\tau\downarrow0$ with $m$ fixed.  Next let $m\to\infty$.
The hypothesis $h_{\mathrm{top}}(F|_Z)=0$ implies
\[
  \limsup_{m\to\infty}\frac{1}{m}\log P(m)=0,
\]
because every fixed-scale entropy rate is nonnegative and bounded
above by the topological entropy.  The right-hand side therefore
tends to zero.  This holds for every $\varepsilon>0$, proving the lemma.
\end{proof}

It remains to compute the entropy on the invariant sets from
Section~\ref{sec:orbits}.  The register map has a possibly
nonmonotone radial cutoff far from the origin, but its first image
lies in a disk where the radius map is nondecreasing.  The following
lemma includes the initial step on the whole sphere.

\begin{lemma}[Entropy of the register dynamics]\label{lemma:register-entropy}
Let $\Sigma=\mathbb C\cup\{\infty\}$ be the sphere, let $A>0$, and
write $D_A=\{v\in\mathbb C:|v|\leq A\}$.  Suppose that a continuous
map $T:\Sigma\to D_A$ satisfies
\[
  T(re^{i\theta})=R(r)e^{i\theta}
  \qquad(0\leq r\leq A),
\]
where $R:[0,A]\to[0,A]$ is continuous and nondecreasing, with
$R(0)=0$.  Then the coordinatewise map
$T^{\times q}:\Sigma^q\to\Sigma^q$ has zero topological entropy
for every integer $q\geq1$.
\end{lemma}

\begin{proof}
Fix a compatible metric $d_\Sigma$ on $\Sigma$ and use its maximum
product metric on $\Sigma^q$.  Fix $\varepsilon>0$.  Uniform continuity
of the polar map
\[
  [0,A]\times\mathbb S^1\longrightarrow D_A,
  \qquad(r,u)\longmapsto ru,
\]
allows a partition of $[0,A]$ into $H$ sufficiently short ordered
intervals and a partition of $\mathbb S^1$ into $K$ sufficiently
short arcs, such that two points have $d_\Sigma$-distance less than
$\varepsilon$ whenever their radii lie in the same interval and their
angular coordinates lie in the same arc.
Use half-open intervals and arcs, with endpoints assigned once,
so that the labels are single-valued.  Let
$b:[0,A]\to\{1,\ldots,H\}$ be the resulting nondecreasing radius
label.

For a radius $r\in[0,A]$ and an integer $N\geq1$, consider its
label vector
\[
  \bigl(b(r),b(R(r)),\ldots,b(R^{N-1}(r))\bigr).
\]
Every coordinate of this vector is nondecreasing as a function of
the initial radius $r$.  Its distinct values therefore form a chain
in the coordinatewise order on $\{1,\ldots,H\}^N$.  Along such a
chain the sum of the coordinates strictly increases whenever the
vector changes, and this sum lies between $N$ and $NH$.  Hence there
are at most
\begin{equation}\label{eq:radius-itinerary-count}
  1+N(H-1)
\end{equation}
radius itineraries.  This orders itineraries by their initial
radii; no monotonicity in time is being asserted.

Each point $v\in D_A$ can be written $v=ru$ with $u\in\mathbb S^1$;
choose $u=1$ when $v=0$.  Its iterates have the representations
$T^jv=R^j(r)u$ with that same $u$, including after a radius becomes
zero.  Thus one angular label, together with the radius itinerary,
gives at most $K[1+N(H-1)]$ descriptions of length-$N$ orbits in
$D_A$.  Points with the same description stay within
$\varepsilon$ throughout those $N$ iterates.

Finally cover the whole compact space $\Sigma^q$ by $C$ sets of
diameter less than $\varepsilon$, and assign each initial point one
cover label.  For $n\geq2$, encode time zero by this label and
encode each coordinate of its first image in $D_A^q$ for the
remaining $N=n-1$ times by the preceding construction.  Points
sharing a complete description stay within $\varepsilon$ in the
maximum product metric for all $n$ times.  Consequently
\begin{equation}\label{eq:register-separated-count}
  s_{\Sigma^q}(n,\varepsilon)
  \leq C\bigl\{K[1+(n-1)(H-1)]\bigr\}^{q}.
\end{equation}
The finite initial cover accounts for all preimages of the disk,
including $\infty$.  The bound grows polynomially in $n$, so its
exponential rate is zero at every $\varepsilon>0$.
\end{proof}

\begin{proposition}\label{prop:zero-entropy}
The map $f$ defined in~\eqref{eq:map} has
$h_{\mathrm{top}}(f)=0$.
\end{proposition}

\begin{proof}
Recall the compact forward-invariant sets of
Proposition~\ref{prop:attraction}:
\[
  \begin{aligned}
    Z_0&=\{(t,z,v_1,\ldots,v_q):t=50\pmod{100},\ z=0\},\\
    Z_\infty&=\{(t,z,v_1,\ldots,v_q):z=\infty,\ v_1=\cdots=v_q=0\}.
  \end{aligned}
\]
On $Z_0$, the clock and main coordinate are fixed and the registers
evolve independently by $T(v)=LS(v)$.  By the choice of $S$ in
Section~\ref{sec:construction}, $T(\Sigma)\subseteq D_{3L}$, and on
that disk its radius map is
\[
  R(r)=Lp(r),\qquad 0\leq r\leq 3L.
\]
This map is continuous, nondecreasing, fixes zero, and takes values
in $[0,3L]$.  Lemma~\ref{lemma:register-entropy} therefore gives
$h_{\mathrm{top}}(f|_{Z_0})=0$.

On $Z_\infty$ the only evolving coordinate is the clock, which
advances by $1$ modulo $100$.  Thus $(f|_{Z_\infty})^{100}$ is the
identity.  For $n\geq100$ the orbit metric $d_n$ on this set equals
$d_{100}$, so its separated-set cardinalities are bounded
independently of $n$ at every fixed scale.  In particular,
$h_{\mathrm{top}}(f|_{Z_\infty})=0$.

For $Z=Z_0\cup Z_\infty$, splitting a separated set between these
two invariant subsets gives
\[
  s_Z(n,\varepsilon)\leq
  s_{Z_0}(n,\varepsilon)+s_{Z_\infty}(n,\varepsilon).
\]
Both terms have zero exponential rate, hence
$h_{\mathrm{top}}(f|_Z)=0$.  Proposition~\ref{prop:attraction}
states that every orbit of $f$ approaches this exact set $Z$.
Lemma~\ref{lemma:localization} now proves
$h_{\mathrm{top}}(f)=0$.
\end{proof}
